\section{Spatial white noise: energy and explicit moment bounds}
\label{S:intrinsic-white-evaluation}

Take $d = \alpha = 1$ and $\Gamma = c\delta_0$, $c \ge 0$, with temporal kernel
$|s-t|^{-a}$, $0 < a < 1$, and data $(1,0)$. The parameter $c$ is covariance
mass; the coefficient of unit noise would be $\sqrt c$. No additional temporal
normalization is inserted. The wave kinetic coefficient is one. For real unit $f
\in H^1(\R)$ put $A = \int_\R f^4$ and $b = \int_\R(f')^2$. The physical
covariance is $A_{\mathrm{sp}}(f) = cA$.

The next calculation recovers the sharp one-dimensional interpolation
inequality recorded by Sz.-Nagy~\cite{sz.nagy:source1941:integralungleichungen}*{(2),
p.~64, and (4), p.~65}, now commonly called a Gagliardo--Nirenberg inequality,
and its hyperbolic-secant optimizer. We give the proof to fix the normalization
of the spatial energy.

\begin{lemma}[A square identity for the distribution function]\label{L:white-cdf-square}
  Every such $f$ has $A < \infty$. With $\mathsf{F}(x) =
  \int_{-\infty}^xf(y)^2\ud y$, for every real $z$,
  \begin{equation}\label{E:white-cdf-square}
    0\le\int_\R[f'-zf(1-2\mathsf{F})]^2
    =b-2zA+z^2/3 .
  \end{equation}
\end{lemma}

\begin{proof}
  On a bounded interval the weak derivative $f'$ is $L^1$ by Cauchy--Schwarz.
  Its integral primitive has weak derivative $f'$ by Fubini. Subtracting it from
  $f$ gives a function of weak derivative zero, hence constant almost
  everywhere: test against derivatives of compact smooth primitives of all
  zero-integral test functions. The resulting locally absolutely continuous
  representatives patch on overlaps, giving one continuous representative. Thus,
  $\rho = f^2$ is locally absolutely continuous with $\rho' = 2ff' \in L^1(\R)$,
  while $\rho \ge 0$ and $\int\rho = 1$. Integrability of $\rho'$ makes $\rho$
  Cauchy on each infinite tail, so it has finite limits at both ends. Its $L^1$
  integrability forces both limits to be zero. It is therefore bounded, and $A =
  \int\rho^2 \le \|\rho\|_\infty\int\rho < \infty$.

  Now $0 \le \mathsf{F} \le 1$, $\mathsf{F}' = \rho$ continuously, and its
  limits are zero and one. The derivative of $\rho(1-2\mathsf{F})$ is
  $\rho'(1-2\mathsf{F})-2\rho^2 \in L^1$. Its endpoint limits vanish.
  Integrating on $[-R,R]$ and sending $R \to \infty$ gives
  \[
    \int_\R f'f(1-2\mathsf{F})=A .
  \]
  The polynomial $P(y) = y-2y^2+(4/3)y^3$ has derivative $(1-2y)^2$. The chain
  rule and the endpoints of $\mathsf{F}$ give
  \[
    \int_\R f^2(1-2\mathsf{F})^2=P(1)-P(0)=1/3 .
  \]
  Absolute integrability holds since this integrand is bounded by $f^2$. For
  each real $z$, both terms in the square in~\eqref{E:white-cdf-square} are in
  $L^2$, and their cross product is absolutely integrable. Expansion and the two
  preceding identities prove the formula. Neither positivity of $f$ nor strict
  increase of its CDF is required. This proves Lemma~\ref{L:white-cdf-square}.
\end{proof}

\begin{proposition}[Exact white spatial energy]\label{P:white-spatial-energy}
  For $c > 0$,
  \[
    R_{\mathrm{sp}}=c/\sqrt3,\qquad
    \mathcal{E}_{\mathrm{sp}}=c^2/6.
  \]
  Both values are zero at $c = 0$.
\end{proposition}

\begin{proof}
  Substitution $z = 3A$ into~\eqref{E:white-cdf-square} gives $3A^2 \le b$. This
  is the exact premise and conclusion of the existing conditional WhiteCDF
  square deduction. Here $A > 0$, since $A = 0$ would imply $f = 0$ a.e. Also $b
  > 0$: zero derivative makes its continuous representative constant and hence
  zero in $L^2$, contradicting unit mass. Thus, $A/\sqrt b \le 1/\sqrt3$.

  For $\ell > 0$ use the real smooth $H^1$ profile $f_\ell(x) =
  \sqrt{\ell/2}\operatorname{sech}(\ell x)$, with $f_\ell' = -\ell\tanh(\ell
  x)f_\ell$. Set $t = \tanh(\ell x)$, so $\ud t = \ell\operatorname{sech}^2(\ell
  x)\ud x$. Direct integration gives
  \[
    \int f_\ell^2=\frac12\int_{-1}^1\ud t=1,\quad
    \int f_\ell^4=\frac\ell4\int_{-1}^1(1-t^2)\ud t=\ell/3,\quad
    \int(f_\ell')^2=\frac{\ell^2}2\int_{-1}^1t^2\ud t=\ell^2/3.
  \]
  The ratio is attained, proving $R_{\mathrm{sp}} = c/\sqrt3$. The physical
  Dirac measure introduces no Fourier constant or nonintegrable kernel
  $|x|^{-1}$. The existing conditional energy deduction applies to $3A^2 \le b$:
  \[
    cA-b/2\le cA-\tfrac32A^2
    =c^2/6-\tfrac32(A-c/3)^2\le c^2/6 .
  \]
  At $\ell = c > 0$ the profile attains $c^2/6$. These are full $H^1$ energy
  profiles; compact approximation equality is
  Proposition~\ref{P:spatial-four-linear}, not a claim that the noncompact
  profile itself is a compact stochastic control. At $c = 0$ the ratio is zero
  and the energy is $\sup(-b/2) = 0$, approached by $f_\ell$ as
  $\ell\downarrow0$. No zero-scale profile is used as a unit vector. The CDF,
  integral and attainment arguments remain analytic; the two finite conditional
  deductions do not formalize them. This proves
  Proposition~\ref{P:white-spatial-energy}.
\end{proof}

\begin{theorem}[Explicit white moment bounds]\label{T:white-moment-bounds}
  For the same original point limit,
  \begin{equation}\label{E:white-moment-bounds}
    \frac23\sqrt{\frac{c}{M_a}}\le C_2(c,a)
    \le\frac23\sqrt{cK_0(a)},\qquad
    K_0(a)=\frac{2^a}{1-a},\quad
    M_a=\frac{\Gamma_E((1+a)/2)}
    {\Gamma_E(1+a)\Gamma_E((1-a)/2)} .
  \end{equation}
  Its time exponent is $(3-a)/2$. These are bounds, not a numerical evaluation
  of the intrinsic supremum.
\end{theorem}

\begin{proof}
  Here $q = 3,k = 2,m = 3/2$. Formula~\eqref{E:spatial-static-energy-score}
  gives
  \[
    F_0(z)=2\left[\frac{(2z)^{1/2}}{8\,3^{3/2}}\right]^{1/2}.
  \]
  For $c > 0$, $\sqrt{2\mathcal{E}_{\mathrm{sp}}} = c/\sqrt3$, hence
  $F_0(\mathcal{E}_{\mathrm{sp}}) = \sqrt c/(3\sqrt2)$ and
  $2^mF_0(\mathcal{E}_{\mathrm{sp}}) = (2/3)\sqrt c$. Substitution into
  Corollary~\ref{C:intrinsic-spatial-bracket}
  proves~\eqref{E:white-moment-bounds} in the original units. All Gamma factors
  are positive finite for $0 < a < 1$. At $c = 0$, $u = 1$ and $C_2 = 0$ by
  Theorem~\ref{T:intrinsic-second-moment}; both displayed bounds vanish. This
  proves Theorem~\ref{T:white-moment-bounds}.
\end{proof}

\begin{corollary}[Uniform relative width of the bound formulas]
  \label{C:white-relative-width}
  For $c > 0$, the ratio of the upper and lower formulas is
  \[
    \sqrt{K_0(a)M_a}
    \le\left[\frac{\sqrt\pi}{2\Gamma_E(5/4)^2}\right]^{1/2},
  \]
  with equality at $a = 1/2$.
\end{corollary}

\begin{proof}
  Duplication and recurrence from Lemma~\ref{L:equilibrium-gamma-beta} give
  \[
    \Gamma_E(1+a)
    =\frac{2^a\Gamma_E((1+a)/2)\Gamma_E(1+a/2)}{\sqrt\pi},
    \qquad
    (1-a)\Gamma_E((1-a)/2)=2\Gamma_E((3-a)/2).
  \]
  Therefore,
  \[
    K_0(a)M_a
    =\frac{\sqrt\pi}{2\Gamma_E(1+a/2)\Gamma_E((3-a)/2)} .
  \]
  Put $x = 1+a/2$, $y = (3-a)/2$, so $x,y > 0$ and $x+y = 5/2$. In Euler's Gamma
  integral the $L^2(0,\infty)$ functions $t^{(x-1)/2}e^{-t/2}$ and
  $t^{(y-1)/2}e^{-t/2}$ have squared norms $\Gamma_E(x),\Gamma_E(y)$ and inner
  product $\Gamma_E(5/4)$. Cauchy--Schwarz, positivity, inversion and the
  increasing square root prove the bound. At $a = 1/2$, $x = y = 5/4$, giving
  equality. Replacing $a$ by $1-a$ exchanges these two arguments, so only this
  ratio of bound formulas is symmetric, not the moment constant or exponent.
  This proves Corollary~\ref{C:white-relative-width}.
\end{proof}

This example retains the all-model theorem rather than replacing its scope. The
formal static exponent is $3/2$ as $a\downarrow0$; no interchange of long-time
and temporal-endpoint limits is asserted. Jensen holds for the same canonical
solution, while $S_1^{\mathrm{pos}}$ remains auxiliary. The read-only Skorohod
comparison at strict spatial Riesz exponent below one does not cover this
spatial-white endpoint. A separate white-time problem cannot be substituted for
$0 < a < 1$. No new Skorohod result, source correction or certified decimal
enclosure is claimed.
